\section*{Introduction}

Energy harvesting refers to the process of capturing and converting ambient energy 
from the surrounding environment into electrical power. Among us, various 
available energy sources are existing, such as thermal, kinetic, solar, and 
radiofrequency (RF). RF energy harvesting has emerged as a particularly promising 
solution for powering autonomous, low-power electronic devices without the need for 
conventional batteries or wired power supplies.

RF energy is available almost everywhere in the electromagnetic spectrum through 
everyday wireless technologies such as Wi-Fi routers, cellular base stations, 
television broadcast towers, and Bluetooth devices. Although the power density of 
these ambient RF signals is relatively low, recent advances in rectenna (rectifying 
antenna) design and ultra-low-power electronics have made it feasible to harvest 
sufficient energy for more and more applications, including wireless sensor networks, 
implantable medical devices, and Internet of Things (IoT) sensors.

A rectenna is the core component of any RF energy harvesting system. It consists of 
a receiving antenna coupled to a rectifying circuit, which converts the incoming 
alternating current (AC) RF signal into a usable direct current (DC) voltage. As 
illustrated in Figure~\ref{fig:rf_system}, a complete RF energy harvesting chain 
typically includes a matching circuit, a voltage multiplier (rectifier), an energy 
storage element, and a microcontroller unit driving the target application. The 
overall performance of such a system is primarily characterized by its RF-to-DC 
power conversion efficiency (PCE), defined as:

\begin{equation}
    \eta_{RF-DC} = \frac{P_{DC}}{P_{RF}} = \frac{V_{DC}^2 / R_L}{P_{RF}}
\end{equation}

\noindent where $P_{RF}$ is the available input RF power, $V_{DC}$ is the output DC 
voltage, and $R_L$ is the load resistance. Maximizing this efficiency requires 
careful co-design of the antenna, the impedance matching network, and the rectifier 
topology.

\begin{figure}[h]
    \centering
    \includegraphics[width=0.45\textwidth]{systeme_rf.png}
    \caption{Block diagram of a complete RF energy harvesting system, from the 
    receiving antenna ($P_{RF}$) to the application unit ($P_{DC}$).}
    \label{fig:rf_system}
\end{figure}

The rectifier topology plays a critical role in determining both the conversion 
efficiency and the output voltage level. Figure~\ref{fig:topologies} presents the 
three conventional rectenna topologies considered in the literature: the series 
topology (a), the shunt topology (b), and the single voltage doubler topology (c). 
In this study, the one-diode (series or shunt) and the two-diode (voltage doubler) 
configurations are designed, simulated, and experimentally compared.

\begin{figure}[h]
    \centering
    \includegraphics[width=0.45\textwidth]{topologies.png}
    \caption{Conventional rectenna topologies: (a) series; (b) shunt; 
    (c) single voltage doubler.}
    \label{fig:topologies}
\end{figure}

The objective of this practical work is to design, simulate, and characterize two AC to DC converter topologies, a one-diode topology and a 
two-diode topology, for low-power RF energy harvesting. Analysis are 
carried out using the Keysight Advanced Design System (ADS) professional software, 
with operating frequencies selected within the 1.5--3~GHz range, covering widely 
used ISM bands such as 2.45~GHz. Both lumped-element (LC) and distributed impedance (microstrip line) matching strategies are investigated and compared.

The report is structured as follows. Section~\ref{sec:task2} presents the choice 
of application, operating frequency, and Schottky diode. Section~\ref{sec:task3} 
details the design and simulation of the one-diode rectifier topology, including 
impedance matching optimization. Section~\ref{sec:task4} extends this analysis to 
the two-diode configuration. Section~\ref{sec:task5} describes the PCB design, 
fabrication process, and experimental measurements. Finally, a comparative 
discussion and conclusion are provided, assessing the suitability of each topology 
for the targeted low-power application.